\begin{table}[H]
\centering
\caption{Summary Statistics}
\begin{threeparttable}
\footnotesize
\begin{tabular}{l ccc ccc ccc}
\toprule
 & \multicolumn{3}{c}{All Mothers} & \multicolumn{3}{c}{DREM=1 Child} & \multicolumn{3}{c}{DREM=0 Child} \\
\cmidrule(lr){2-4} \cmidrule(lr){5-7} \cmidrule(lr){8-10}
 & Mean & SD & N & Mean & SD & N & Mean & SD & N \\
\midrule
Employed & 0.698 & 0.459 & & 0.593 & 0.491 & & 0.705 & 0.456 & \\
Hours/week & 27.5 & 18.6 & & 23.9 & 19.6 & & 27.8 & 18.5 & \\
Annual wages (\$) & 28,287 & 40,934 & & 19,967 & 32,047 & & 28,811 & 41,374 & \\
In labor force & 0.746 & 0.435 & & 0.666 & 0.472 & & 0.752 & 0.432 & \\
Age & 40.0 & 7.1 & & 39.4 & 7.2 & & 40.0 & 7.1 & \\
College degree & 0.329 & & & 0.212 & & & 0.337 & & \\
Married & 0.710 & & & 0.567 & & & 0.719 & & \\
White & 0.714 & & & 0.719 & & & 0.713 & & \\
\midrule
Observations & & & 2,400,329 & & & 140,835 & & & 2,259,494 \\
\bottomrule
\end{tabular}
\begin{tablenotes}[flushleft]
\small
\item \textit{Notes:} ACS 1-Year PUMS, 2008--2019. Sample: women aged 25--54 who are the reference person or spouse in households with children aged 5--17. DREM=1 indicates the household has at least one child with cognitive difficulty. All statistics weighted using ACS person weights (PWGTP). Wages are nominal annual wages/salary income. Employment defined as ESR $\in \{1,2,4,5\}$.
\end{tablenotes}
\end{threeparttable}
\label{tab:summary}
\end{table}

